\section{Smooth passages and spectral densities}\label{sec:passages}

A volume-preserving coordinate system, also called a chart below, means
a smooth volume-preserving identification of the physical torus with $\T^3$.
All coordinate systems in the construction preserve normalized volume.
An intermediate transformation has the form
\begin{equation}\label{eq:twist}
 S(\theta,u,a)=(\theta+\phi(u),u,a),
 \qquad (\theta,u,a)\in\T^3,
\end{equation}
where $\phi\in C^\infty(\T;\R)$ has $\phi'>0$ on an open arc
$I$. We call $I$ a good arc. Derivatives on an arc are understood in an
increasing real lift of that arc.

A packet is a function $A(u,a)\ee(m\theta)$, where
$m\in\Z\setminus\{0\}$ and $A$ is smooth with compact support in
$I\times\T$. A finite packet list is disjoint if, for any two entries with
the same $m$, their coefficient supports are disjoint. Entries with
different $m$ need no support restriction. Packets may be grouped into
colors; we write $F_r$ for the sum of color $r$. Disjointness implies
\begin{equation}\label{eq:packet-orthogonality}
 \langle U_S^n F_r,U_S^k F_l\rangle=0
 \quad(r\ne l,\ n,k\in\Z).
\end{equation}
Indeed $U_S^n(A\ee(m\theta))=\ee(nm\phi(u))A\ee(m\theta)$,
and either angular orthogonality or disjoint coefficient supports applies.
The spectral measure of one packet is
\begin{equation}\label{eq:packet-spectral-measure}
 (m\phi\bmod1)_*\bigl(|A(u,a)|^2\dd u\dd a\bigr).
\end{equation}
It has a smooth density because its coefficient support is compactly
inside the good arc. We use the inner product linear in its first
variable, so a vector of density $p$ has
$\langle U_S^nF,F\rangle=\int_\T\ee(nz)p(z)\dd z$.
A measurable spectral multiplier $v(z)$ acts on its
coefficient by multiplication by $v(m\phi(u)\bmod1)$, whenever the result
belongs to $L^2$. These statements follow directly from the displayed
formula for $U_S^n$, first for trigonometric polynomials and then by
$L^2$ approximation. The pushforward formula, absolute continuity,
and multiplier rule also apply to nonsmooth coefficients in $L^2$
supported in the good arc; their densities need not be smooth.

For later comparisons, if two vectors have spectral densities $p_F,p_G$
for the same transformation, then
\begin{equation}\label{eq:density-norm-bound}
 \|p_F-p_G\|_1
 \le (\|F\|_2+\|G\|_2)\|F-G\|_2.
\end{equation}
To see this, test their spectral measures against an arbitrary function
of absolute value at most one and use the contraction property of its
spectral multiplier. The total variation norm here and below is the
$L^1$ norm of densities, without a factor $1/2$.

\begin{lemma}[Passage]\label{lem:passage}
Fix \eqref{eq:twist}, its physical volume-preserving coordinate chart,
and a finite disjoint packet list. Choose a positive integer $c$ so large
that the compact ranges $m(c+\phi(\T))$, for the distinct old indices $m$,
are pairwise disjoint and avoid zero. Set $V=c+\phi$.
Fix $J\in C^\infty(\T;\R)$ with $J'>0$ on an open arc $I'$ and fix
$0\le\psi\le1$ in $C_c^\infty(I')$.
For integers $N\to\infty$, choose $\Delta_N$ so that
\begin{equation}\label{eq:passage-rounding}
 |\Delta_N|\le N^{-1},\qquad
 k_N:=N\left(\int_\T V(u)\dd u+\Delta_N\right)\in\Z,
 \qquad H_N'=N(V+\Delta_N).
\end{equation}
Put $\sigma_N(a)=N^{-1}+N^{-2}J(a)$. Then
\begin{equation}\label{eq:passage-map}
 T_N(\theta,u,a)=
 \bigl(\theta+H_N(u+\sigma_N(a))-H_N(u),\,
 u+\sigma_N(a),\,a\bigr)
\end{equation}
is a volume-preserving smooth torus diffeomorphism, and both $T_N$ and
$T_N^{-1}$ converge smoothly to $S$ and $S^{-1}$, also in the fixed
physical coordinates. In the new coordinates
\begin{equation}\label{eq:passage-chart}
 (\theta',u',a')=C_N(\theta,u,a)
             =(u,a,\theta-H_N(u)),
\end{equation}
it is the twist of speed $\sigma_N$ and good arc $I'$.

Here are its packet and density properties. For an old packet indexed
by $\alpha$, with index $m_\alpha$ and amplitude $A_\alpha$, define
\begin{equation}\label{eq:stationary-image}
 E_\alpha=\{(m_\alpha V(u),a):(u,a)\in\supp A_\alpha\}.
\end{equation}
One may choose pairwise disjoint, arbitrarily small fixed neighborhoods
of these compact sets, and smooth cutoffs $\chi_\alpha(s,a)$ supported
in those neighborhoods and equal to one near $E_\alpha$. Retain the
finitely many new packets
\begin{align}
 G_{\alpha,j}(\theta',u',a')
   &=\ee(j\theta')\ee(m_\alpha a')
       \psi(u')\chi_\alpha(j/N,u')C_{\alpha,j}(u'),
       \label{eq:retained-packets}\\
 C_{\alpha,j}(a)
   &=\int_\T A_\alpha(u,a)\ee(m_\alpha H_N(u)-ju)\dd u.
       \label{eq:exact-passage-coefficient}
\end{align}
All their angular indices are nonzero for sufficiently large $N$.
For each $\alpha$,
\begin{equation}\label{eq:transport-norm}
 \left\|\sum_jG_{\alpha,j}\circ C_N
       -\psi(a)A_\alpha(u,a)\ee(m_\alpha\theta)\right\|_2=o(1).
\end{equation}
They are disjoint as packets, and their moduli are independent of the
new spare coordinate $a'$. Every constant attached to an old packet is
preserved by this linear transport and these cutoffs.

For a single old packet write, with zero extension outside its branch,
\begin{equation}\label{eq:W-definition}
 W(s,a)=\frac{|A(u,a)|^2}{|m\phi'(u)|},\qquad s=mV(u).
\end{equation}
Let $K(s,y)$ be the density of the pushforward of
$W(s,a)\psi(a)^2\dd a$ under $a\mapsto sJ(a)$. It is continuous and
compactly supported in $(s,y)$, with its $s$ support away from zero.
For any fixed positive integer $P$ and residue $h\pmod P$, the unfolded
temporal density of the retained packets with $j\equiv h\pmod P$ differs
in $L^1(\R)$ by $o(1)$ from
\begin{equation}\label{eq:fine-density-unfolded}
 Q_{N,h}(t)=\sum_{k\in\Z}K(t,Nt-h-kP).
\end{equation}
Consequently their density on $\T$ differs in $L^1(\T)$ by $o(1)$ from
$\sum_{l\in\Z}Q_{N,h}(z+l)$. The assertions hold after summing over
any of the finitely many colors. All $o(1)$ terms refer to the old data,
$c,J,\psi,P$, and support neighborhoods fixed before $N$ tends to
infinity. In particular, $N$ may be required to exceed any prescribed
bound.
\end{lemma}

\begin{proof}
\emph{Geometry and smooth convergence.}
Choose $k_N$ by rounding $N\int V$ to an integer and integrate
\eqref{eq:passage-rounding}. Then $H_N(u+1)-H_N(u)=k_N$, so the map
$u\mapsto H_N(u)\bmod1$ is globally defined. The map $C_N$ is a smooth
torus diffeomorphism of Jacobian one, with inverse
$(\theta',u',a')\mapsto(a'+H_N(\theta'),\theta',u')$.
Direct substitution gives
\[
 C_NT_NC_N^{-1}(\theta',u',a')
       =(\theta'+\sigma_N(u'),u',a').
\]
This proves the global assertions, except convergence. The integral
identity
\[
 H_N(u+\sigma_N(a))-H_N(u)
       =N\int_0^{\sigma_N(a)}(V(u+t)+\Delta_N)\dd t
       =V(u)+O_{C^k}(N^{-1})
\]
holds for each fixed $k$, with constants depending on the fixed $c,J$
and $\phi$. Since $c$ is an integer, addition of $V$ and of $\phi$ agree
on $\T$. The inverse of \eqref{eq:passage-map} is
\[
 (\theta,u,a)\longmapsto
 (\theta-H_N(u)+H_N(u-\sigma_N(a)),u-\sigma_N(a),a),
\]
and the same estimate proves inverse convergence. Composing with a
fixed physical chart and its inverse preserves convergence in every
fixed smooth norm. No bound uniform over future coordinate charts is
being asserted.

\emph{Fourier coefficients.}
Work on a compact
subarc of $I$ containing all old coefficient supports in its interior.
For $s=j/N$ the phase in \eqref{eq:exact-passage-coefficient}, divided
by $N$, has critical equation
\begin{equation}\label{eq:critical-equation}
 m(V(u)+\Delta_N)=s
\end{equation}
and second derivative $m\phi'(u)$. It has at most one critical point
on this arc. With $u_*=u_*(s-m\Delta_N)$ on that branch, set
\[
 L_{m,j}(a)=
 \frac{A(u_*,a)}{\sqrt{N|m\phi'(u_*)|}}
 \ee\left(mH_N(u_*)-ju_*+\frac{\operatorname{sgn}(m)}8\right),
\]
and set this to zero wherever the amplitude at the critical point is
zero or there is no point on the branch. Uniformly for $j/N$ in any
fixed compact interval and $a\in\T$,
\begin{equation}\label{eq:stationary-uniform}
 C_{m,j}(a)-L_{m,j}(a)=o(N^{-1/2}).
\end{equation}

For general stationary-phase background see \cite{Hormander2003};
we give the compact-support one-dimensional argument needed here.
Enlarge the support arc twice inside $I$. Critical points
near the original support then range over a fixed compact interval on
which the absolute second derivative is bounded below. A smooth change
of variable, depending smoothly on the critical point, reduces the
divided phase to its critical value plus
$\operatorname{sgn}(m)y^2/2$; its Jacobian at zero is
$|m\phi'(u_*)|^{-1/2}$. After putting $v=\sqrt N y$, the integral on
$|v|\le D$ converges uniformly to the constant amplitude times the
quadratic exponential integral there. The remaining local tails are
$O(N^{-1/2}/D)$, uniformly in the parameters, by integrating against
$\ee(\operatorname{sgn}(m)v^2/2)$: the rescaled amplitudes have uniformly
bounded suprema and total variation. The part away from the critical
point is $O(N^{-1})$ by integration by parts. Finally
\[
 \lim_{D\to\infty}\int_{-D}^D\ee(\pm v^2/2)\dd v
          =\exp(\pm i\pi/4);
\]
this follows by inserting a Gaussian factor, evaluating the Gaussian
integral, and removing the factor using the same tail bound.
Taking first $N\to\infty$, then $D\to\infty$, proves
\eqref{eq:stationary-uniform}. When the potential critical point is
outside the enlarged support interval, the integral is uniformly
nonstationary. This also covers zero amplitudes and the ends of the
branch.

There are $O(N)$ indices in a fixed bounded $j/N$ range. Outside a
sufficiently large such range, repeated integration by parts gives,
for any fixed positive integer $q$,
$|C_{m,j}(a)|\le C_q|j|^{-q}$ uniformly in $a$ and $N$.
Indeed the phase derivative has size comparable to $|j|$, and its
higher derivatives are $O(N)$ with $N=O(|j|)$ there. Thus
\eqref{eq:stationary-uniform} implies
\[
 \sum_j\int_\T |C_{m,j}-L_{m,j}|^2\dd a=o(1).
\]
The $L^2$ norm of the leading array is bounded by a Riemann sum.
Cauchy--Schwarz therefore gives the useful mass estimate
\begin{equation}\label{eq:summed-stationary-mass}
 \sum_j\int_\T
 \left||C_{m,j}(a)|^2
       -\frac1N W(j/N-m\Delta_N,a)\right|\dd a=o(1).
\end{equation}

\emph{Cutoffs and transport.}
The map $(u,a)\mapsto(mV(u),a)$ is injective on $I\times\T$.
It takes disjoint compact supports with the same $m$ to disjoint
compact sets. Our choice of $c$ separates the images for different
$m$. The stated cutoffs $\chi_\alpha$ therefore exist. For large $N$
they equal one on the support of the corresponding leading array,
since its $s$ coordinate is shifted by only $m_\alpha\Delta_N$.
Because each cutoff is one on the shifted leading-array support, the
deleted exact coefficient mass is bounded by the summed squared
$C_{m,j}-L_{m,j}$ error, which tends to zero by the preceding estimate
and Parseval's identity. The same argument is valid after
multiplication by $\psi$, proving \eqref{eq:transport-norm} and
\eqref{eq:summed-stationary-mass} for the retained arrays with the
factor $\psi^2$ on the right. The cutoffs make separation exact, not
merely asymptotic. Different old $m$ never contribute to the same
retained $j$, and at a shared $j$ the remaining supports in $u'=a$
are disjoint. The factor $\ee(ma')$ in
\eqref{eq:retained-packets} proves the asserted independence of the
new spare-coordinate modulus, also after combining packets at one
index. Linearity proves the statement about attached constants.

\emph{Fine temporal density.}
We first remove rounding before any fine-scale sampling. Smoothness and
compact support of the fixed $W$ give
\begin{equation}\label{eq:rounding-mass-removal}
 \frac1N\sum_j\int_\T
 |W(j/N-m\Delta_N,a)-W(j/N,a)|\dd a=o(1).
\end{equation}
For example, a uniform bound on $\partial_s W$ and $O(N)$ nonzero
summands give an $O(|m|/N)$ bound. The same holds with $\psi^2$ and
for any residue subset. Pushforward contracts total variation, so
\eqref{eq:summed-stationary-mass} and
\eqref{eq:rounding-mass-removal} allow us to replace the retained
squared coefficient measures by
$N^{-1}W(j/N,a)\psi(a)^2\dd a$, at total error $o(1)$, before
pushing to temporal frequency. In particular, the actual lattice
$j/N$ is never translated by $m\Delta_N$.

On the retained arc, the change of variable $y=sJ(a)$ gives
\begin{equation}\label{eq:K-jacobian}
 K(s,y)=\frac{W(s,a)\psi(a)^2}{|s|J'(a)},
 \qquad a=J|_{I'}^{-1}(y/s),
\end{equation}
with zero extension. The compact support of $\psi$ inside $I'$ makes
this extension continuous (in fact smooth away from $s=0$).
The real temporal frequency of a new angular index $j$ is
\begin{equation}\label{eq:output-frequency}
 t=j\sigma_N(a)=j/N+(j/N)J(a)/N.
\end{equation}
The density of the pushforward of
$N^{-1}W(j/N,a)\psi(a)^2\dd a$ under this map is exactly
$K(j/N,Nt-j)$. The factor $N$ in the density change of variable
cancels the coefficient mass $N^{-1}$.

Choose $Y$ so that $K(s,y)=0$ for $|y|>Y$. For fixed $t$, at most
$2Y+2$ integers contribute to either
$\sum_j K(j/N,Nt-j)$ or $\sum_jK(t,Nt-j)$; on every contributing
term $|j/N-t|\le Y/N$. Uniform continuity of $K$, and a fixed
bounded interval containing the $t$ supports for large $N$, prove
that replacing $j/N$ by $t$ costs $o(1)$ in $L^1(\R)$.
The same proof works after restriction to $j\equiv h\pmod P$.
This is \eqref{eq:fine-density-unfolded}. Folding modulo one
contracts the $L^1$ norm and finishes the proof.
\end{proof}

It is useful to emphasize the normalization in
\eqref{eq:fine-density-unfolded}. The periodization
$\sum_k K(t,y-kP)$ is a density relative to length on the circle of
length $P$, so its mean value is
\begin{equation}\label{eq:residue-mean}
 \frac1P\int_\T W(t,a)\psi(a)^2\dd a.
\end{equation}
For $P=1$, folding $t=z+l$ leaves the fast phase $Nz$ unchanged,
since $N$ is an integer. For general $P$, the folded expression has
phase $N(z+l)-h$ modulo $P$; no divisibility condition on $N$ is
imposed.

\begin{lemma}[Fast-variable averaging]\label{lem:fast-average}
Let $P>0$, and let $G(t,y)$ be continuous, compactly supported in $t$,
and $P$-periodic in $y$. Then, uniformly in real offsets $\beta_N$,
\begin{equation}\label{eq:fast-average}
 \int_\R G(t,Nt+\beta_N)\dd t
 \longrightarrow \frac1P\int_\R\int_0^P G(t,y)\dd y\dd t.
\end{equation}
In particular, suppose $w\ge0$ is continuous and compactly supported
in a compact set $K\subset\R\setminus\{0\}$. For continuous
integrable densities $\rho,\widetilde\rho$ define
\[
 \mathcal P_{s,P}\rho(y)
     =\sum_{k\in\Z}|s|^{-1}\rho((y+kP)/s).
\]
Assume these periodizations are jointly continuous for $s\in K$.
Then
\begin{align}\label{eq:law-diagonal-comparison}
 &\lim_{N\to\infty}
 \int_\R w(s)
 |\mathcal P_{s,P}\rho(Ns+\beta_N)
       -\mathcal P_{s,P}\widetilde\rho(Ns+\beta_N)|\dd s\\
 &\hspace{12mm}=
 \frac1P\int_\R w(s)\int_0^P
 |\mathcal P_{s,P}\rho(y)-\mathcal P_{s,P}\widetilde\rho(y)|
       \dd y\dd s
 \le\frac{\|w\|_1}{P}\|\rho-\widetilde\rho\|_1.\notag
\end{align}
Joint continuity is automatic, for example, if the densities are
continuous with bounds $C/(1+y^2)$, or are continuous and compactly
supported. Neither density need have mass one.
\end{lemma}

\begin{proof}
Partition the $t$ axis into intervals of length $P/N$ on each of
which $Nt+\beta_N$ traverses one period. On each interval replace
the slow argument $t$ of $G$ by an endpoint. Uniform continuity on
the fixed compact support bounds the integrated error by a constant
times its modulus of continuity at $P/N$, with an additional
$O(N^{-1})$ boundary error. The resulting averages are Riemann sums
for the right side of \eqref{eq:fast-average}. This proves the first
assertion, uniformly in the offset. Apply it to the continuous
absolute difference in \eqref{eq:law-diagonal-comparison}.
The integral over one period of that difference is bounded by
\[
 \sum_k\int_0^P |s|^{-1}
       |\rho((y+kP)/s)-\widetilde\rho((y+kP)/s)|\dd y
       =\|\rho-\widetilde\rho\|_1.
\]
Finally, on a compact set of nonzero $s$, quadratic tail bounds give
a uniform summable bound $C'/(1+k^2)$ on the periodizing summands,
for $0\le y\le P$. Uniform convergence proves joint continuity.
Compactly supported densities are handled in the same way with
only finitely many nonzero summands on that compact set.
\end{proof}

In the special case $|A(u,a)|$ is independent of $a$,
\eqref{eq:W-definition} has the form $W(s,a)=w(s)$.
If $\widetilde\rho\dd y=J_*(\psi^2\dd a)$, then
$K(s,y)=w(s)|s|^{-1}\widetilde\rho(y/s)$. Thus replacing the
realized jitter law by a regular target law $\rho$ in the
fine-density prescription costs, after taking $N$ sufficiently
large, at most
$\|w\|_1\|\rho-\widetilde\rho\|_1+o(1)$ for the full density.
One sums this estimate over the fixed packet list and folds modulo
one. This is an averaged estimate along the fast variable, not a
pointwise estimate of periodized densities from total variation.
It is applied only after \eqref{eq:rounding-mass-removal}, so the
family being averaged is fixed.

\begin{lemma}[Realization of jitter laws]\label{lem:law-approx}
For any strictly positive smooth probability density $\rho$ on
$\R$ and any $\eps>0$, there are a smooth real circle function $J$,
an open arc $I'$ with $|\T\setminus I'|<\eps$ and $J'>0$ on $I'$,
and $0\le\psi\le1$ in $C_c^\infty(I')$, such that
\[
 J_*(\psi^2\dd a)=\widetilde\rho(y)\dd y,
 \qquad \widetilde\rho\in C_c^\infty(\R),
 \qquad \|\rho-\widetilde\rho\|_1<\eps.
\]
Here the cutoff law is a subprobability measure; no renormalization
is made. Also $\int_\T(1-\psi^2)\dd a<\eps$.
\end{lemma}

\begin{proof}
Let $F(y)=\int_{-\infty}^y\rho(v)\dd v$. Strict positivity makes
$F$ a smooth diffeomorphism from $\R$ to $(0,1)$. Choose
$0<\delta<\min(1/8,\eps/4)$ and a smooth function $\xi$ compactly
supported in $(\delta/2,1-\delta/2)$ and equal to one on
$[\delta,1-\delta]$. In the circle coordinate $0<a<1$, set
$J(a)=\xi(a)F^{-1}(a)$ and extend it by zero near the identified
endpoints. It is a smooth bounded circle function, and on
$I'=(\delta,1-\delta)$ it has derivative
$1/\rho(F^{-1}(a))>0$. Choose $\psi\in C_c^\infty(I')$ with
$0\le\psi\le1$ and $\psi=1$ on $[2\delta,1-2\delta]$.
Change of variables gives
\[
 \widetilde\rho(y)=\rho(y)\psi(F(y))^2\le\rho(y),\qquad
 \|\rho-\widetilde\rho\|_1
       =1-\int_\T\psi^2\dd a\le4\delta<\eps.
\]
Its support is compact because that of $\psi$ is compactly inside
$(0,1)$. All the claims follow.
\end{proof}

\begin{lemma}[Weighted neutral refresh]\label{lem:neutral}
Given a twist, its fixed physical chart, a finite disjoint smooth
packet list with color densities $p_r$, a positive integer $P$,
and a positive tolerance, one can perform a passage with the
following simultaneous properties. Its physical map and inverse
lie in any prescribed smooth neighborhood of the old ones. Its
new good-arc complement is arbitrarily small. Every color vector
is transported within the prescribed $L^2$ tolerance, retaining
constant tags and disjointness. The new packet moduli are
independent of the new spare coordinate. The output densities
$p'_r$ and the densities $p'_{r,h}$ contributed by output indices
$j\equiv h\pmod P$ satisfy
\begin{equation}\label{eq:neutral-densities}
 \|p'_r-p_r\|_1<\eps,
 \qquad \|p'_{r,h}-p_r/P\|_1<\eps
 \quad(0\le h<P),
\end{equation}
where $\eps>0$ may be prescribed arbitrarily.
The integer $c$ may first be required to exceed any fixed bound,
and the retained supports may use any sufficiently small fixed
enlargements of \eqref{eq:stationary-image}. If colors were
disjoint in their old angular indices alone, their output angular
index sets are still disjoint. The final passage integer $N$ may
exceed any prescribed bound.
\end{lemma}

\begin{proof}
Fix the old finite data and choose $c$ as in Lemma~\ref{lem:passage},
also satisfying its prescribed lower bound. Identify all but a
small interval of the $a$ circle with an interval in $\R$.
Choose an open arc $I'$ whose complement is as small as requested,
and $0\le\psi\le1$ compactly supported there and equal to one
except on a set of sufficiently small length. The boundedness of
the finitely many old amplitudes allows us to make both their
$L^2$ cutoff errors and the removed squared masses arbitrarily
small. Fix this cutoff before proceeding.

For a positive parameter $L$, take a smooth periodic function
$J_L$ equal to $La+\beta$ on $I'$, extending it smoothly across
the complementary interval. Such an extension exists by taking
a smooth cutoff equal to one on the closure of $I'$ in a slightly
larger real coordinate interval. In particular $J_L'=L$ on $I'$.
Put $f_s(a)=W(s,a)\psi(a)^2$, extended by zero in this real
coordinate. The periodization of its pushforward by $sJ_L$ is
\begin{equation}\label{eq:neutral-riemann-sum}
 \sum_{k\in\Z}\frac1{|s|L}
 f_s\left(\frac{(y+kP)/s-\beta}{L}\right).
\end{equation}
This is $1/P$ times a Riemann sum for $\int f_s(a)\dd a$, of mesh
$P/(|s|L)$, with an arbitrary lattice offset and possibly reversed
orientation. The error is bounded by
\begin{equation}\label{eq:neutral-riemann-error}
 \frac{\|\partial_a f_s\|_{L^1(\R)}}{|s|L}.
\end{equation}
Indeed compare each lattice value times the mesh with the integral
on its cell, using the fundamental theorem of calculus and summing
the absolute derivative. Since $s$ ranges over a fixed compact set
away from zero, this bound tends to zero uniformly in $s$ and $y$.
It remains valid for arbitrary fixed dependence of the original
amplitude on $a$.

Choose $L$ large enough for these uniform errors to be as small as
needed for the finite packet list. Now $J_L$ and all its norms are
fixed. Choose the support cutoffs in Lemma~\ref{lem:passage}, and
then take $N$ sufficiently large. Its fine-density rule, together
with \eqref{eq:neutral-riemann-sum}, approximates each unfolded
residue density by $P^{-1}\int W(s,a)\psi(a)^2\dd a$ in $L^1$.
Folding modulo one gives $p_r/P$, up to the small cutoff loss;
the integer $mc$ in $mV$ has no effect on this folding. The
estimate is uniform in the periodic phase, so the folding offsets
$Nl$ cause no difficulty. Summing the residue estimates gives the
total density estimate, upon initially making each error smaller
by the fixed factor $P$ if necessary. Norm transport, exact
support properties, and smooth proximity follow from the same
passage lemma. All constraints on $N$ are lower bounds imposed
after the physical chart, the amplitudes, $c$, the cutoff, $L$,
and the support neighborhoods have been fixed.
\end{proof}

\begin{remark}\label{rem:finite-family-passages}
A common passage can handle finitely many packet lists, requiring
disjointness only within each list. Choose $c$ using the union of
their old indices, choose the neutral cutoff and slope to meet all
their finitely many mass and density bounds, and choose $N$ above
all the resulting lower bounds. Support cutoffs are chosen
separately for each list, disjoint within that list. The norm and
density conclusions then hold for every list; no orthogonality
between different lists is asserted.
\end{remark}

All these operations preserve real data. Include together the
packets $(m,A)$ and $(-m,\overline A)$, attach conjugate tags, and
choose $\chi_{-m}(s,a)=\chi_m(-s,a)$. The exact coefficients obey
$C_{-m,-j}=\overline{C_{m,j}}$, and all other cutoffs are real.
Thus conjugate packets and conjugate colors can be retained at
every passage.
